\documentclass[output=paper,colorlinks,citecolor=brown]{langscibook}
\ChapterDOI{10.5281/zenodo.18845602}
\author{Roberta Medda-Windischer\orcid{}\affiliation{Institute for Minority Rights, Eurac Research (Bolzano, Italy)} and Mattia Zeba\orcid{/0000-0001-7609-1234}\affiliation{Institute for Minority Rights, Eurac Research (Bolzano, Italy)}}
\title[Linguistic diversity of new minorities in Europe]{Linguistic diversity of new minorities in Europe: A crossroads of rights and duties}
\abstract{Increasing international mobility flows have presented both opportunities and challenges for social cohesion and integration in Europe. Language, as a tool for communication and a marker of identity, plays a pivotal role in ensuring equality of access to education, the labor market, and justice. However, many European states face difficulties in balancing the need for integration with the recognition of linguistic diversity, particularly concerning new minorities arising from migration flows. These challenges have underlined the urgent need for effective governance of linguistic diversity.

This paper explores the governance of linguistic diversity in the context of integrating new minorities in Europe. It argues for expanding the scope of minority rights — traditionally reserved for historical minorities — to include new minorities. This approach addresses gaps in international frameworks, particularly regarding identity and language rights, and fosters a more inclusive diversity governance model. Recognizing and protecting the linguistic rights of new minorities is essential to promoting cohesion and belonging, reducing the risk of alienation and conflict, and avoiding forced assimilation. The paper highlights the dual need for supporting minority languages while maintaining proficiency in a common public language to ensure state functionality and societal integration.

Using an interdisciplinary methodology, the paper combines legal and linguistic approaches to analyse language policies for new minorities. It examines three key areas — education, the labor market, and the judicial system — demonstrating how language rights intersect with societal inclusion and equity, and offering recommendations for balancing diversity and unity in modern Europe.}

%move the following commands to the "local..." files of the master project when integrating this chapter

% \usepackage{tabularx}
% \usepackage{langsci-optional}
% \usepackage{langsci-gb4e}
% \bibliography{localbibliography}
% \newcommand{\orcid}[1]{}

\IfFileExists{../localcommands.tex}{
   \addbibresource{../localbibliography.bib}
   \usepackage{tabularx,multicol}
\usepackage{url}
\urlstyle{same}
\usepackage{textcomp}
 
\usepackage{langsci-optional}
\usepackage{langsci-lgr}
\usepackage{langsci-gb4e}
\usepackage{langsci-branding}

% \usepackage{fontspec}
\babelprovide{georgian} % Georgian language
\babelprovide{armenian} % Armenian language
\babelfont[georgian,armenian]{rm}{DejaVu Sans}
\newfontfamily\georgianfont{DejaVu Sans}
\newfontfamily\armenianfont{DejaVu Sans}

\usepackage{dialogue}
\usepackage{attrib}
\usepackage{attrib} % speaker labels in Koca's chapter
\usepackage{phoenician} % alf, bet in Mengozzi's chapter

   

\makeatletter
\let\thetitle\@title
\let\theauthor\@author
\makeatother

\newcommand{\togglepaper}[1][0]{
   \bibliography{../localbibliography}
   \papernote{\scriptsize\normalfont
     \theauthor.
     \titleTemp.
In Saloumeh Gholami \& Ulrich Mehlem (eds.), \textit{Identity formation, language and migration dynamics: Minorities in the MENA region and the German diaspora}, Berlin: Language Science Press [preliminary page numbering]
   }
   \pagenumbering{roman}
   \setcounter{chapter}{#1}
   \addtocounter{chapter}{-1}
}

\newbool{bookcompile}
\booltrue{bookcompile}
\newcommand{\bookorchapter}[2]{\ifbool{bookcompile}{#1}{#2}}


\newcommand{\georgian}[1]{{\georgianfont #1}}
\newcommand{\armenian}[1]{{\armenianfont #1}}

\newcommand{\textarab}[1]{{\arabicfont \beginR #1} \endR}
\newcommand{\texthebrew}[1]{{\hebrewfont \beginR #1} \endR}
\newcommand{\textsyriac}[1]{{\syriacfont \beginR #1} \endR}


\newfontfamily\phoenicianfont{Tuffy Regular.ttf}

   %% hyphenation points for line breaks
%% Normally, automatic hyphenation in LaTeX is very good
%% If a word is mis-hyphenated, add it to this file
%%
%% add information to TeX file before \begin{document} with:
%% %% hyphenation points for line breaks
%% Normally, automatic hyphenation in LaTeX is very good
%% If a word is mis-hyphenated, add it to this file
%%
%% add information to TeX file before \begin{document} with:
%% %% hyphenation points for line breaks
%% Normally, automatic hyphenation in LaTeX is very good
%% If a word is mis-hyphenated, add it to this file
%%
%% add information to TeX file before \begin{document} with:
%% \include{localhyphenation}
\hyphenation{
    par-a-digm
    Ver-to-vec
    Go-go-lin
    sa-gen
    mo-misch
    ant-wor-te-te
    dix-wa-zim
    spre-chen
    B-Verf-GE
    Swe-dish
    re-mains
    schrei-ben
    di-pey-vin
    Chris-ten-sen
    Wes-sen-dorf
    Ja-va-khi-shvi-li
    ana-lysis
}

\hyphenation{
    par-a-digm
    Ver-to-vec
    Go-go-lin
    sa-gen
    mo-misch
    ant-wor-te-te
    dix-wa-zim
    spre-chen
    B-Verf-GE
    Swe-dish
    re-mains
    schrei-ben
    di-pey-vin
    Chris-ten-sen
    Wes-sen-dorf
    Ja-va-khi-shvi-li
    ana-lysis
}

\hyphenation{
    par-a-digm
    Ver-to-vec
    Go-go-lin
    sa-gen
    mo-misch
    ant-wor-te-te
    dix-wa-zim
    spre-chen
    B-Verf-GE
    Swe-dish
    re-mains
    schrei-ben
    di-pey-vin
    Chris-ten-sen
    Wes-sen-dorf
    Ja-va-khi-shvi-li
    ana-lysis
}

   \boolfalse{bookcompile}
   \togglepaper[23]%%chapternumber
}{}

\begin{document}

\maketitle

\section{Introduction}\label{ch2.4.ss.1}

In most European countries, diversity governance has emerged as a complex and challenging issue, currently holding a prominent position in political and public discourse. This is largely due to the increasing number of peoples—especially migrants and asylum seekers—with distinctive identities in terms of language, culture, or religion in urban as well as in more peripheral and rural contexts with varying degrees of permanence. While economic stakeholders and decision-makers generally acknowledge the valuable contributions to the labor force and positive effects on demographic structures, the presence of large immigrant and refugee communities presents complex challenges in the sphere of integration with regard to cultural and linguistic differences, protection of individual and group rights, and the preservation of social cohesion and unity. Identity claims, including linguistic claims, and economic rights however are not separate issues, as diversity claims very often arise precisely in labor and employment environments. Accommodation of diversity and labor market inclusion are indeed complementary and intimately intertwined, thus one cannot be solved without the other.
\largerpage[2]

The right to identity in diversity is also intimately linked to language.\footnote{The \citeauthor{scc1988_ford_quebec} has clarified this point in \textit{Ford v. Quebec} (\citeyear{scc1988_ford_quebec}) by saying that `[l]anguage is not merely a means or medium of expression; […]. It is a means by which a people may express its cultural identity. It is also the means by which one expresses one's personal identity and sense of individuality.' } Language is not only an important tool to communicate, but also for identity formation; language rights are thus based on a `duality of structures' \citep{giddens1984} in the form of the \textit{multidimensionality} and \textit{multifunctionality}: languages as an \textit{instrumental medium} of communication as well as a substantial element for \textit{personal and social identity formation} in the construction of `meaning, trust and order' \citep{eisenstadt1995}.\footnote{See also \citet[34]{VertovecWessendorf2005}.} In this regard, language has been referred to as a gate and as a tie \citep{kraus_kazlauskaite_gurbuz2014}.\footnote{\citeauthor{kraus_kazlauskaite_gurbuz2014} relate this approach to Ralf Dahrendorf's distinction between options and ligatures, according to which ``Options are possibilities of choice''; they provide us with ``structural opportunities for choice, thereby offering a template for our individual choices and decisions''. Ligatures, in contrast ``are allegiances; one might call them bonds or linkages as well''. \citet{dahrendorf1979}, cited by \citet[521]{kraus_kazlauskaite_gurbuz2014}.}

Given the relevance of both instrumental and identity-related interests in language, the \textit{politics} of language have often represented an ideological, political, and legal arena for linguistic, sociocultural, and political control based on processes of language standardization through selection of a particular language, usually that of the most powerful group, discouraging at the same time the use of other languages or even varieties (dialects) of the same language in the public sphere, thereby encouraging users to develop loyalty and pride in it \citep{marko_constantin2019}.\footnote{\citeauthor{sonntag_cardinal2015} refer to `language regimes' as concerning not only language practices but also ``conceptions of language and language use as projected through state policies and as acted upon by language users''. \citep[6]{sonntag_cardinal2015}.} Indeed, in linguistically diverse settings, individuals face the necessity of choosing an interaction language, often leading to an asymmetrical preference for one language over others. When this process becomes systematic, resembling situations where certain groups consistently yield, it mirrors scenarios of social imbalance \citep{van_parijs2011}. From a multicultural perspective, ``linguistic justice'' would involve safeguarding the linguistic rights of minority groups, thereby ensuring their ability to use their language in public and preventing the injustice and inequality that would arise if they were forced to shift to another language \citep{deschutter2007,van_parijs2011,mowbray2012,alcalde2015,riera-gil2019,gazzola_wickstrom_fettes2023,medda-windischer_constantin2024}.

Language standardization \citep{costa_de_korne_lane2017,zeba2020} and the adoption of monolingual/territorial approaches \citep{barry2001,van_parijs2011} --- with or without legally declaring the majority language the official language of the state --- go hand in hand with the need for standardized public education and literacy demands for the labor force which is thus postulated a `normal,' instrumental requirement. As a consequence, bi- or multilingual education and the use of languages other than the dominant or official language seem to require an exemption from this rule of monolingualism in public education, the economy, administration or politics \citep{marko_constantin2019}. 

\textit{Language policies} have thus political and social consequences; they can reinforce or diffuse conflicts, tensions or social unrest between language groups, they accelerate language loss or facilitate language revitalization, and they can be instruments of inclusion or exclusion \citep{sonntag_cardinal2015}. Restrictions on language use and barriers in the acquisition of language proficiency may exclude individuals from education, political participation, and access to justice, thus underscoring the critical role of language in societal structures \citep{osce_hcnm2017,un2004,skutnabb-kangas2008}. 

Against this backdrop, this chapter seeks to address the following research questions: How can the demands for linguistic diversity and political unity be reconciled to create a cohesive and stable political community that addresses the legitimate aspirations of minorities, including new minorities?\footnote{On social cohesion, see the seminal study by \citet{putnam2001,stolle_soroka_johnston2008,kelser_kraaykamp2010,portes_vickstrom2011}.} What public policies should be implemented to achieve this balance between diversity and unity? How can the risk of essentializing minority cultures through political mobilization, which might lead to divisive us-versus-them antagonisms, be mitigated?\footnote{On essentialisation of minorities, see \citet{marko_medda-windischer2018,vertovec_wessendorf2005}.} What strategies can support multiple integration into the country of residence while allowing individuals to maintain identification with the culture of their country of origin or minority group? Further, how can institutional segregation from common institutions or shared societal sectors—particularly the public education system and the labor market—be avoided, thereby preventing socioeconomic downward assimilation into ethnic ghettos and the complete cultural marginalization of minority groups? And finally, how can expanding the scope of minority rights, traditionally associated with historical or \textit{old} minorities, to include \textit{new} minorities originating from migration, contribute to more effective diversity governance and address the gaps in international frameworks, particularly regarding identity and language rights? How might such an expansion foster inclusion and a sense of belonging among minority groups, while preventing forced assimilation and mitigating the risks of resistance and alienation?

Indeed, if it is certain from a human and minority rights perspective that \textit{fundamental human rights and liberties} such as the right to existence, identity in diversity, freedom of expression and association must be accorded to all human beings, it is less certain which are the \textit{specific legal obligations of states} in language matters towards persons belonging to historical minorities and new minority groups stemming from migration. Moreover, these obligations are less clear when they are connected to claims vis-à-vis the state to adopt special measures to ensure appropriate conditions for the preservation and development of group identities. 

In the following, this chapter examines language rights and duties for new minorities arising from recent migration flows within Europe by adopting an interdisciplinary approach that integrates legal and linguistic studies. The analysis draws on legal and policy documents, existing literature, and empirical research conducted in this field. The chapter is structured in three sections: the first section provides the basis for the following parts by analyzing the diversity of old and new minorities and their alleged dichotomy; the second section explores the normative framework pertaining to language rights and duties for new minorities, in particular as residence and citizenship are concerned; in the third part three major ambits in which language policies for new minorities have a particular relevance, namely education, the labor market and the law enforcement system, are analyzed. These ambits have been selected on various grounds: education serves as a foundational domain where linguistic skills are cultivated, cultural understanding is fostered, and opportunities for upward mobility are shaped. The labor market relies heavily on linguistic proficiency as a determinant of employability, workplace communication, and professional advancement, making language a key factor in economic integration. Lastly, the judicial system highlights the importance of language in ensuring equitable access to justice, comprehension of legal rights and obligations, and the ability to participate fully in all aspects of societal life. The chapter concludes with some final remarks concerning future challenges and possible solutions.

Prior to analyzing the various sections comprising the chapter, it is important to provide some terminological clarifications. The linguistic debate has distinguished between `migrant languages' and `immigrant languages.' The former are characterized by a high degree of individual mobility and of local fluctuations in the number of groups, so that a migrant language cannot take root in a given territory. The latter are instead characterized by low degrees of mobility and fluctuations, so that they can influence the ``local idiomatic order'' \citep{bagna_machetti_vedovelli2003}. In other words, the term `immigrant languages' identify those instances where an immigrant community shows ``a long-lasting, if not definitive, migratory project'' \citep{orioles2007} in a given territory, as well as a manifested will to ``maintain their original identity, culture and language'' \citep{orioles2007,panzeri2016}. Other authors \citep{trifonas_aravossitas2018} instead have used terms such as, among others, `languages of origin' \citep{makarova2014}, `ethnic languages' \citep{saint-jacques1979}, `community languages' \citep{wiley2005}, `home languages' \citep{yeung_marsh_suliman2000} or simply `mother tongues.' There is still a debate on the correct use of these terms, as well as on the need to substitute `mother tongue' with a more scientific term such as `first language.' In particular, the CoE has been increasingly using `first language' to refer to ``the language acquired first in early childhood, which has a special status for the child as it is acquired as part of its discovery of the world in primary socialisation and is in fact crucial to the success of the child in developing normally'' \citep{coe2007}.  Instead, `mother tongue' ``includes emotional associations of ideas with `family' and `origins' which `first language' excludes'' \citep{coe2007}. Specifically in the field of education, the term `heritage languages' has been gaining consensus since the 1970s to define in general ``language spoken by the children of immigrants or by those who immigrated to a country when young'' \citep{cho_shin_krashen2004}, and in particular ``the incompletely learned home language arising from the phenomenon of language shift and the switch to the dominant language that is characteristic in the case of immigrants and their descendants'' \citep{trifonas_aravossitas2018,polinsky_kagan2007}. Because of this process, heritage language speakers are defined as ``[...] individuals who have been exposed to an immigrant or minority language since childhood, who are also very competent in the majority language spoken in the more widespread linguistic community'' \citep{montrul2015}. Therefore, such a term not only indicates a language spoken by immigrants but also hints at its `health condition,' that is, the level of competence in which it is spoken. 

\section{Old and New Minorities: Setting the Grounds }\label{ch2.4.ss.2}

In almost all European countries live communities whose members have a distinct language, religion, or socio-cultural practices as compared to the rest of the population and who have become minorities through the redrawing of international borders, having seen the sovereignty of their territories shift from one country to another. These are ethnic groups that have not achieved statehood on their own, for various reasons, and that have now become part of a larger country (or several countries): they are the so-called `old minorities' \citep{brubaker2004,smith2010,keating2002,mabry_et_al2013,marko_constantin2019}. In many, but not all cases, their co-ethnics may be numerically or politically dominant in another state, which they therefore regard as their `external national homeland' or kin-state \citep{osce_hcnm2008}. Similar to the case of old minorities, in most European countries there are groups formed by individuals and families who have left their original homelands to emigrate to other countries: these are the so-called `new minorities.'\footnote{It is important to note that `new minorities,' as any other groupings of people, are communities based on fictitious constructions and narratives, or, following the classic Andersonian definition of national communities, they are  ``imagined communities'' resulting from the ``social construction of reality'' \citep{Anderson1983}. New minorities then come into existence through processes of group formation in terms of social organisation such as family reunification, economic networks and forming civil and political associations. Commenting on Brubaker's definition of ethnic groups according to which ``[ethnic groups] are collective cultural representation […] that sustain the vision and division of the special world in racial, ethnic, or national terms'' \citep{brubaker2004}, Marko clarifies that ``(w)henever people define this situation as `real' the consequences following from their actions are no less `real' than the `existence' of things (Marko 2018).} In most cases, their reasons are economic, although political reasons play an increasingly important role as well. New minorities thus consist of migrants and refugees, and their descendants, who are living in a country other than that of their origin, on a basis that is more than merely transitional. They encompass multiple variations of situations ranging from acute refugee situations to economic migrants and their family reunification.

It must be acknowledged that the term `new minorities' is subject to difficulties and criticism because it encompasses an enormous variation of situations. The use of the term is intended to refer to `distinct' groups, and it by no means implies a weakening of their status. It aims to offer additional legal tools with which to respond to their specific needs for protection \citep{acfc2016,acfc2001a}. Moreover, the term `new minorities' is broader than the term `migrants,' as it encompasses not only the first generation of migrants, but also their descendants, extending to include second and third generations of individuals with a background of migration, many of whom have been born in the country of immigration and who cannot objectively or subjectively be subsumed under the category of `migrants.' Finally, the term `new minorities' emphasizes the diversity of the individuals concerned, as well as their related individual and collective rights, whereas the term `migrants' does not \citep{medda-windischer2017}.

Drawing from Laponce's dated distinction between `minorities by force' and `minorities by will' \citep{laponce1960}, Walzer regards immigrants as having chosen to leave their original cultures, with the awareness that the success of their decisions depends upon integrating into the mainstream of their new societies \citep{walzer1995}. In such cases, ethnic diversity arises from the \textit{voluntary} decisions of individuals or families to uproot themselves and join another society. In contrast, Walzer argues, old minorities remain focused on their historic homelands. These groups find themselves in a minority position, not because they have uprooted themselves from their homeland, but because their homeland has been incorporated within the boundaries of a larger state. In most cases, this incorporation was \textit{involuntary}, resulting from conquest, colonization or the ceding of territory from one imperial power to another. Under these circumstances, it is argued, minorities are rarely satisfied with models based on non-discrimination, individual rights and, eventually, integration. As proposed by Walzer, they desire `national liberation' (i.e. some form of collective self-government) in order to ensure the continued development of their distinct culture. This differentiation could be called into question, however, largely because it is debatable whether migrants have truly made a voluntary `choice' to migrate. This applies not only to refugees or those fleeing from wars or natural disasters, but also to `labor migrants' escaping from economic distress.

Also, against an expansion of language rights to new minorities, \citet{kymlicka1996} argues that if migration is voluntary, this affects the claims immigrants can make in the receiving society, limiting them to polyethnic rights safeguarding cultural and religious practices, not extending to language rights. He acknowledges the difficulty in drawing a line between involuntary refugees and voluntary immigrants, proposing that the only long-term solution is addressing the unjust international distribution of resources, as treating economic refugees as national minorities in a new country may not provide an appropriate redress for injustices that must ultimately be resolved in the original homeland. Similarly, \citet{patten2006} contends that requiring immigrants to relinquish their claim on official status for their languages is a necessary condition to secure interests in democratic government and language maintenance. He argues that a commitment by public institutions to officially recognize every immigrant language, along with the languages of established groups, would be crippling to democratic self-government. Therefore, established community members (members of 'national groups') have priority over immigrants in claiming official language rights precisely because they are established. This distinction is evident both at the local and supranational levels, such as within the EU, where national languages, with English increasingly prominent, are celebrated most, regional or minority languages less, and immigrant languages least \citep{extra2017}.

A crucial issue in discussing language rights is that claims of minorities—old and new minorities alike—are often perceived as a challenge and antagonistic to the traditional model of culturally homogeneous nation states because both groups seek to increase, within this model, opportunities to express their identities and diversities at individual and group level. Moreover, old minorities and new groups stemming from migration are often perceived as `foreigners' in the state where they live. Members of historical and new minorities, as long as they are not absorbed into the national body through assimilation, are considered – even after generations - as loyal to their kin-state or to the original state whose citizens the first generation migrants had been.

Historically, new minorities stemming from migration have reacted very differently than historical minorities to societies dominated by an ethnic majority. Unlike historical minorities whose cultural traditions may predate the establishment of the state that their members are now citizens of, in general few migrant groups object to the requirement that they must learn the official language of the country of residence as a condition of citizenship \citep{van_oers2014}, or that their children must learn the official language in school. Migrants usually accept that their life chances and those of their children depend largely on the participation in mainstream institutions operating in the majority language \citep{kymlicka2001}. 

Problems related to integration of second and third generations can be quite acute with regard to new minorities, though this also applies to a certain extent to traditional minorities, especially in the case of mixed marriages. The children of second and third generations are in fact subjected to the decisions taken by their parents and their living with two cultures and languages can be perceived either as an enriching experience or, more frequently, as an excessive burden. This is due to the fact that often the second and third generations of new minorities' descendants have less cultural distance from the society of the country of residence than their parents, but they have not reached a satisfactory degree of integration from a socio-economic viewpoint. In comparison to their parents, their expectations are as high as those of their `pairs' without a migration background, and consequently, the risk of frustration and alienation from mainstream society may be very high \citep{alba_holdaway2013,wihtol_de_wenden2015}.

In the discussion on minorities and diversity, an ongoing debate is whether the scope of application of international treaties pertaining to minorities that are usually applied to historical, old minorities can be extended to new minority groups stemming from migration \citep{hofmann2006_impact_of_international_norms,kymlicka2007,medda-windischer2017,medda-windischer2009,joppke2010,jackson_preece2010}.\footnote{Specifically on language rights for old and new minorities, see, \citet{poggeschi2013}.} The positions in this regard are extremely diversified: among states, some have adopted rather narrow views firmly opposing the extension of minority provisions to new minorities \citep{coe1995,coe1997}, others have instead pragmatically applied some provisions to new groups \citep{acfc2001b,acfc2004},\footnote{In this regard, it has to be noted the proactive position of the Czech Republic that has officially recognized a new minority group as a national minority, namely the Vietnamese community, the third largest minority group in the Czech Republic. With the new status, Czech-Vietnamese receive support from the government's budget to preserve Vietnamese culture, tradition and language. In particular, the official minority status ensures them the right to use Vietnamese language at offices as well as courts, set up Vietnamese broadcast or TV programs, among others; see \citet{kascian_vasilevich2013,trlifajova2013}.} while others have taken a more blurred position. 

Most international bodies dealing with minorities have adopted in this regard an inclusive approach, especially the Advisory Committee on the Framework Convention \citep{acfc2016,acfc2001c,acfc2002a,acfc2002b}, the European Commission for Democracy Through Law (2007), the UN Human Rights Committee,\footnote{The UN Human Rights Committee clarified that the terms of Article 27 ICCPR indicate that individuals designated to be protected need not be citizens of a state party, and that a ``State party may not therefore restrict the rights under Article 27 to its citizens alone''. The UN HR Committee continues: ``Just as they need not be nationals or citizens, they need not be permanent residents. Thus, migrant workers or even visitors in a State party constituting such minorities are entitled not to be denied the exercise of those rights''. See \citet{UNHRC1994}. See also \citet{UN HRC1981}.} the UN Working Group on Minorities \citep{eide2000}, and the OSCE High Commissioner on National Minorities which seems to have extended its mandate to new minority groups stemming from migration \citep{osce2004,osce_hcnm2006,ekeus2006,osce_hcnm2012}.

The approach taken by the Advisory Committee on the Framework Convention with regard to new minority groups originating from immigration is particularly relevant in this regard. Due to the significant proportion of non-citizens – including migrant workers – in the total population of some countries, the Advisory Committee found that it would be possible to consider the inclusion of persons belonging to these groups in the application of the Framework Convention on an \textit{article-by-article} basis, and noted that the authorities of the countries concerned should consider this issue in consultation with those concerned.\footnote{For instance, in 2007, the ACFC mentioned Turkish \textit{Gastarbeiter} in Germany as a group that could benefit from certain rights covered by the Framework Convention for Protection of National Minorities (FCNM) (Second Opinion on Germany 2007). One year later, it welcomed the appointment by the Department for Intercultural and Integration Affairs of Vienna of `a person specifically in charge of dealing with problems facing the Roma, whether autochthonous or those with migrant background' \citep{acfc2008}. Moreover, the ACFC pointed out that Austria would need to provide support for the preservation of the cultural and linguistic heritage of Czechs, Slovaks, and Hungarians living in Vienna \citep{acfc2008}. These minority communities have a historical presence in Vienna, but nowadays they include among their number those who have moved to Austria in recent decades.} In particular, the ACFC emphasized that ``inclusive language policies should cater for the needs of everybody based on their different characteristics and needs, including persons belonging to national minorities living outside their traditional areas of settlement, immigrants and non-citizens'' \citep[24]{acfc2016}. 

It has to be acknowledged, however, that, as it will be further discussed in this chapter, any decision to bring minorities with a recent international mobility background within the scope of application of international and/or national instruments pertaining to minorities is bound to be mainly a political one.

\section{Normative Framework on Language Rights and Duties of New Minorities}\label{ch2.4.ss.3}

The discourse around rights and duties in the linguistic sphere of new minorities raises a number of questions, in particular, whether the use of a particular language for either official or non-official purposes is a right or not,\footnote{One vexing, yet unresolved question concerns whether minority rights have a collective or an individual dimension. As stated by Marko, ``[t]hese two forms of rights not only can, but even must be used cumulatively when organising equality on the basis of difference'' (\citealt{marko1997}, see also \citealt{hofmann1998}). Ultimately, the true issue concerns whether the groups that human beings form are free and whether the members of these groups are able to live in dignity, including with regard to the maintenance and development of their identities.} and even if we can speak in terms of `rights,' what pragmatic limitations can legitimately be placed on the exercise of language choice. As seen earlier, rights are balanced against other (often relatively abstract) interests such as general societal interests – e.g. requirements of the `common good' or of `public policy' – or more specific economic or political interests as spelled out as possible ``substantive limitations'' for state interference in the European Convention of Human Rights or the EU Fundamental Rights Charter \citep{shuibhne2002}.

Hence, discussing language rights and duties of new minorities refers to the wider debate on \textit{social integration}, that is the focus of this chapter, and \textit{system integration} developed and conceptualized in scholarly literature \citep{marko_constantin2019,gordon1964,berry1997_again,esser2001_integration}.

\textit{Language rights and duties} of new minorities from international, European and national legislation can be classified around three categories: 1) Language rights necessary to preserve and maintain cultural ties with the country of origin, or that of the parents or minority group in case of second/third generations, and, 2) Language duties (or requirements) which are necessary to be able to communicate and participate in the overall society, and 3) State obligations to protect and to promote both the preservation of the language of origin as well as the learning of the language of the country of residence which is not in every situation identical with the rights and duties under 1) and 2). Just to illustrate this with a not only hypothetical example: Shall parents have a right to home-schooling of their children (under 1) so that children (under the age of 14) would have no right to attend a public school with mandatory learning of the state language as language of instruction?

\subsection{Language Rights of New Minorities}\label{ch.2.4.ss.3.1}

International norms dealing with language rights of new minorities are those foreseen, for instance, in the UN Convention of the Rights of the Child that extends the right to education far beyond equality of access to education and includes provisions concerning the development of the child's cultural identity, language and values of the child's country of origin.\footnote{See, art. 29 (1) (c) of the UN CRC 1990.} 

Most international instruments for the protection of migrants, such as the United Nations' 1990 International Convention on the Protection of the Rights of All Migrants Workers and Members of Their Families (2003),\footnotemark{} the Council of Europe`s 1977 Convention on the Legal Status of Migrant Workers (1983) or the EU Directive on the Status of Third-Country Nationals who are Long Term Residents (LTRD) (2003), refer to the teaching of the migrant workers' mother tongue for their children. These instruments however contain also references to possible requirements of `integration' that include linguistic training in the language of the settlement country.\footnote{See, artt. 5(2) and 15(3) of the EU Council Directive (2004); and also art. 14 of the CoE European \textit{Convention} on the Legal Status of \textit{Migrant Workers} (1977).} Moreover, the aim of the provisions concerning mother tongue tuition is mainly the return of these children to the country of origin of their parents rather than the protection and promotion of their identities.\footnote{The COE Migrant Workers' Convention (Art. 15) reads specifically: ``The Contracting Parties concerned shall take actions by common accord to arrange, so far as practicable, for the migrant worker's children, special courses for the teaching of the migrant worker's mother tongue, to facilitate, inter alia, their return to their State of origin.''. See also artt. 1 and 67 of the UN \textit{International Convention} on the Protection of the Rights of All \textit{Migrant Workers} and Members of Their Families (2003).} The sole EU legislative measure concerning new minorities' language education is Directive 486/EEC, entitled \textit{On the Education of the Children of Migrant Workers}, which was adopted by the then European Community in 1977. The Directive's scope is limited to children of workers who are nationals of other Member States (MSs) and it aims to ease the difficulties associated with the initial reception and integration of such children into the educational school system of the country of residence. The Directive 486/EEC requires EU Members to ensure, among others, the promotion of the teaching of the mother tongue and culture of the country of origin of such children, along with regular education and in cooperation with the country of origin; and ensure free tuition in their territory, adapted to the specific needs of such children, in particular the teaching of an official language of the settlement State. As said, Directive 486/EEC deals only with the education of children who are EU citizens, and it does not address a substantial part of the challenge posed by most recent migration, namely the education of children who are third-country nationals \citep{eu_coe_fra2015}.\footnotetext{The UN Migrant Workers' Convention is considered the most comprehensive and detailed legal treaty addressing the rights of migrant workers, and it is a treaty that can potentially serve as an important tool for strengthening the legal protection of members of new minorities and contributing to their integration into society. However, this treaty is currently in force mainly among countries of emigration, none of which are EU Member States. As of 11 December 2024, the Migrant Workers' Convention has been signed by 40 countries and ratified by 60, with only seven being CoE Member States: Albania, Armenia, Azerbaijan, Bosnia and Herzegovina, Montenegro, Serbia, and Turkey. See \url{https://treaties.un.org/Pages/ViewDetails.aspx?src=TREATY\&mtdsg_no=IV-13\&chapter=4\&clang=_en}.}

In general, the role of the European Union in the field of education is limited to supporting and promoting the educational policies of its Member States. The latter retain full control in establishing their school systems and school curricula, provided that they apply the principle of equality and prohibit discrimination. At the EU level, however, the issue of education of children who are unable to speak the language of the country of residence has increasingly gained increased attention. The European Commission, in particular, facilitates the exchange of good practices on the integration of new minorities and funds relevant projects across the different levels of education.\footnote{For an overview of the activities of the European Commission in the field of education for migrants and refugees settled in EU MSs, see the web page: \url{http://ec.europa.eu/education/policy/migration_en}} This is being attributed to the transformation that the EU has undergone in the last years due to the greater mobility associated with the latest enlargements and the large influx of migrants and asylum seekers from non-EU countries \citep{jacob_owens2022,milano2023}. 

\subsection{Language Duties of New Minorities}\label{ch.2.4.ss.3.2}

Following the rise of migration in Europe, a growing number of European countries has been adopting norms and regulations providing and/or requiring language tests or courses for residence and, sometimes, even admission to residence, as well as for the acquisition of citizenship.\footnote{For a recent comparative overview see \citet{eu_fra2017,van_oers_groenendijk_kostakopoulou2010,iglesias_sanchez2009,jesse2017}.}

In general, new countries of immigration, such as Greece, Ireland, Italy and Spain, are less concerned about proficiency in the national language than countries which have a longer history of immigration, particularly when these countries are former countries of emigration which in a very short lapse of time have become countries of immigration. The approach in these countries is quite the opposite end of that in countries which have experienced a great deal of family-based immigration and have chosen to introduce a number of language requirements for the issue of many documents for migrants and their families.


With regard to EU legislation, the EU Directive on the Status of Third-Country Nationals who are Long Term Residents (LTRD) (2004) provides that third-country nationals are entitled to the long-term resident status after residing legally and continuously for five years in the territory of a Member State before their application for status (LTRD 2004). The status entitles long-term residents to equal treatment with nationals in a number of areas and enhanced protection against expulsion. As seen earlier, the Directive permits states to impose integration requirements as a condition for acquiring long-term resident status in the first place,\footnote{Art. 5(2) LTRD (Conditions for acquiring long-term resident status) provides: ``Member States may require third-country nationals to comply with integration conditions, in accordance with national law''. and Art. 15 LTRD (Conditions for residence in a second Member State) reads: ``Member States may require third-country nationals to comply with integration measures, in accordance with national law. This condition shall not apply where the third-country nationals concerned have been required to comply with integration conditions in order to be granted long-term resident status, in accordance with the provisions of Article 5(2). Without prejudice to the second subparagraph, the persons concerned may be required to attend language courses''.} and to impose language requirements as a condition for access to education and training.\footnote{Art. 11(3)( b) LTRD. } The integration requirements foreseen in the LTRD have a precise antecedent in the citizenship legislation of some MSs, which require `integration' to be understood as a result-oriented process proven through language proficiency and knowledge of the country, for instance history, culture, tradition, legal system, Constitution, geography, etc.) in order to obtain citizenship \citep{baubock_et_al2006,acosta_arcarazo2011}.


As of residence permits and integration requirements, Italy, for instance, has introduced an integration agreement that must be signed by a third-country national entering Italy for the first time. According to the law a language test must also be passed in order to obtain the long-term residence permit (Law n. 94 2011). There are exceptions for those who are under 14 years of age and those who suffer from limitations to learning languages due to their age or physical condition (Italian Immigration Act n. 286 2011).

In other cases, such as the Flemish community in Belgium, newly arrived immigrants are required to take an officially provided language course, though the granting of a residence permit is not linked to the course. This is due to the fact that permission to enter the country and the granting of residence permits is dealt with at federal level, whereas the regional level is responsible for integration policies and their implementation \citep{coe2014}.


According to Bauböck, obligatory language requirements for the right to stay in the country of settlement are legitimate only `when they assume a convergence of public interests and private interests of the immigrants themselves' \citep{baubock2003}. Hence, mandatory language courses for adults would be justifiable `as a form of benign paternalism' \citep{baubock2003}. From this point of view, language courses can secure migrants' long-term interests regarding social upward mobility as opposed to short-term interests in earning income in low-skilled jobs leading to forms of `down-ward assimilation' \citep{marko_constantin2019}. However, the legitimacy of defining such interests `from above,' that is from outside migrant communities, remains questionable \citep{vertovec_wessendorf2005}.


As of the acquisition of citizenship, over the past decade, numerous European countries have increasingly mandated language proficiency tests and examinations assessing knowledge of various aspects of the host country, such as its history and traditions, as an essential prerequisite for obtaining citizenship \citep{acosta_arcarazo2011}. Sometimes, the conditions required in order to acquire a long-term residence permit are the same as those needed to acquire citizenship, but in general there has been a transition by which integration conditions to obtain citizenship have gradually also been adopted for the acquisition of permanent residence \citep{acosta_arcarazo2011}.

For instance, according to the Austrian Nationality Act, it is a precondition for the acquisition of Austrian nationality that the applicant for naturalization demonstrates knowledge of the German language and a basic knowledge of the democratic legal order and the history of Austria and the respective federal state (Austrian Nationality Act 2017). Also in Germany, language is an important requirement for naturalization: the applicant for naturalization has to demonstrate `sufficient' oral and written skills of the German language (German Nationality Act 2016). The `sufficient' nature of the skills can be demonstrated in different ways, such as the participation of the applicant in a language course when applying for a residence permit \citep{guild_groenendijk_carrera2009}. In the UK applicants applying for naturalization must demonstrate that they have sufficient understanding of English (or Welsh or Scottish Gaelic), in addition to sufficient knowledge of life in the UK (UK Nationality, Immigration and Asylum Act 2002).\footnote{Regarding the command of English for speakers of other languages, this implies ``the ability to hold a conversation on an unexpected topic, which is workable, though not perfect, English'' \citep{guild_groenendijk_carrera2009}. The main goal of the test is however not to assess the level of command of the English language but rather the knowledge of the candidate about the British society. In the UK therefore the level of command of the English language is only implicitly tested: the UK assumes that by passing the immigration test on the knowledge of British society the applicant has already demonstrated that his command of the language is sufficient \citep{guild_groenendijk_carrera2009}.}

Exceptions are applied in most countries, generally for mental physical or intellectual diseases or handicaps, or age (for instance, in Austria, UK and Germany) \citep{guild_groenendijk_carrera2009}, but also if the mother tongue of the person is already the language to be tested (Austria and Germany) \citep{guild_groenendijk_carrera2009}.

\section{Main Spheres of Life with Language Policy Implications for New Minorities  }\label{ch2.4.ss.4}

Language, as previously noted, serves as a pivotal marker of identity and plays a critical role in the construction of identity at both individual and collective levels \citep{acfc2012}. Beyond its function as an emblem of cultural and personal identity, language is also an essential instrument for facilitating integration into society. Its significance extends across multiple spheres of life, influencing social, economic, and institutional interactions. When considering the relationship between language and new minorities, three key dimensions stand out as particularly significant. First, education represents a critical sphere where language skills are developed, cultural awareness is nurtured, and pathways to social and economic advancement are established. Second, the labor market places a strong emphasis on linguistic competence, which is essential for securing employment, facilitating effective workplace communication, and enabling career progression, thereby making language a central component of economic integration. Finally, the judicial system underscores the vital role of language in guaranteeing fair access to justice, understanding legal rights and responsibilities, and actively engaging in various aspects of societal life. Together, these dimensions underscore the multifaceted role of language in addressing the challenges and opportunities associated with new minorities within contemporary societies.

\subsection{Education }\label{ch2.4.ss.4.1}

While persons belonging to new minorities shall not be discriminated against in their access to education, and, as previously discussed, the UN Convention on the Rights of the Child explicitly extends the right to education beyond mere equality of access to include provisions for fostering the child's cultural identity, language, and the values of their country of origin \citep{liefaard_sloth-nielsen2017,tobin2019}, a broader and more complex issue arises regarding the extent to which these individuals can assert a right for their identity and culture, including their language, to be meaningfully incorporated into the educational process.

If language plays a role in various aspects of life concerning new minorities, education serves as the foundational aspect of language protection and a substantive guarantee for the continuous use of a language. The lack of education in the minority language jeopardizes its transmission between generations, especially considering social and economic incentives, such as broader opportunities and improved social status, favoring the majority language \citep{deschutter_robichaud2016}. Furthermore, the process of education has a profound impact, positively or negatively, on a young person's sense of identity. Education has in fact been described as a ``powerful instrument for the achievement of social engineering'' \citep{thornberry_estebanez2004}. Education can help to strengthen and further develop that identity while creating awareness and tolerance of other cultures existing in the same society. Educational policies should therefore combine a focus on the universal values, the practical needs of children thereby facilitating their participation in socioeconomic life, and the respect for their distinct cultural traditions and identities \citep{acfc2024}.

The political and legal progression from `language as a problem,' to `language as a right' – and to `language as a resource' \citep{ruiz1984,gogolin1994,ruiz2010,dejong_li_zafar_wu2016} – is a difficult, hotly contested process. Indeed, states seem to see the granting of language rights to minorities a highly divisive issue and ultimately leading to the disintegration of the state \citep{skutnabb-kangas2008,vertovec_wessendorf2005}. As Palermo has rightly pointed out: ``Languages of the minorities are often seen as the main threat to the development of the state language, thus something against which the state language must be protected. This often creates a clash between laws aimed at protecting the minority languages and state language laws'' \citep{palermo2012}.

In education there is thus a constant tension between these aspects: How to preserve, on the one hand, the distinctive identity of a minority and, on the other, how to contribute to social cohesion? How far does the teaching of and education in new minorities' languages lead to the retreat of minority members into their communities or encourages social integration beyond the family?

During the first waves of postwar immigration from postcolonial areas and from southern Europe during the 1970s and 1980s issues of language rights and education for new minorities have been largely dismissed \citep{twitchin_demuth1985,verma_bagley1979,Ager1996LanguagePolicy}). However, proficiency in the language of origin of new minorities is increasingly considered to be of great importance for pupils \citep{eurydice2009}, since it can facilitate the learning process of the state language or language of instruction\footnote{As explained by \citet{medda-windischer2009}: while ``the Interdependence Theory based on the concept of Common Underlying Proficiency (CUP) or the integrated source of thought for both languages accommodates the process of skills and knowledge transfer between the two languages'', ``the Minimal Threshold of Linguistic Competence conditions the positive transfer of skills between languages by suggesting that it is necessary to obtain a certain secure level of proficiency in both languages in order to be experiencing the benefits of bilingualism''. See \citet{cummins1976,cummins1978,cummins1991,cummins2000}. Some authors have criticized such hypotheses, especially the Minimal Threshold theory, also supporting the idea that the positive effect of bilingualism comes from the learning of a second language and not the maintenance of the mother tongue. See, among others, \citet{esser2012}.} and thus stimulate pupils' development and participation in various spheres of life. For some students, knowledge of the language of their parents might open up additional opportunities for their educational and professional development and could improve their chances on the job market, although the evidence supporting this assumption is unclear \citep{pendakur_pendakur2002}. In addition, knowledge of the language of origin of their parents can contribute to secure the self-esteem and identity of children belonging to new minorities, helping them to preserve and intensify their social ties with members of their community in the country of origin and in the settlement country \citep{bankston_zhou1995,eurydice2009,zhang2008}.

Nevertheless, it is also acknowledged that children who do not speak, read, or write the language of instruction to the level of their peers perform less well in school \citep{oecd2001}.\footnote{On educational inequality see, \citet{baysu_devalk2012,heath2008,phalet_heath2011,entorf_minoiu2005,levels_dronkers2008,oecd2006}, and on dropout rates, \citet{kalmijn_kraaykamp2003}.} The Programme for International Student Assessmentt (PISA) has demonstrated that being of immigrant background still constitutes a disadvantage with respect to school success, putting the responsibility on public schools to establish equal chances for every citizen \citep{mehlem_maas_schroeder2004,christensen_stanat2007}. This argument has been defined as the `institutionalized reproduction of inequality' as a factor for new minorities' underachievement \citep{marko2013,radtke_gomolla2002,vertovec_wessendorf2005}.

Helping children of new minorities maintain and develop both languages – the state language or language of instruction and the language of the country of origin of their parents – is therefore a worthwhile, though difficult, goal \citep{christensen_stanat2007}. Regardless of the different attitudes that lead to education models and teaching methods, finding a balance between the three aims of education, namely universal values, practical needs of the child and respect for distinct cultural traditions and identities, is often described, especially by teachers, as difficult or problematic. The problems range from complications of the teaching assignment and decreased educational quality because of restricted linguistic competences of the children, to difficulties to convince parents and administrations that it is profitable to teach languages of countries of origin of new minorities like Turkish or Arabic, which for many seem to be largely irrelevant to European societies \citep{decarlo2005}.

\largerpage[-1]
Mainstream monolingual ideologies have a particularly strong impact on schools, where teachers are often (more or less consciously) highly skeptical  about the advantages of heritage language maintenance. This trend has been confirmed by various studies all across Europe. For instance, in the UK ``there is a lack of support in mainstream education for heritage language, which has made it necessary for parents and families to accept full responsibility'' \citep{weekly2020}; this ``creates an implied language policy, when one actor does not contribute, it is left to other actors to assume responsibility, in this instance parents, families and communities'' \citep{weekly2020}. Similarly, quantitative studies in Greece ``have indicated that a very large number of teachers (48.2\%) consider their pupils' heritage languages a hindrance to the learning of Greek, while another large number stated that the learning of heritage languages should concern the immigrant communities themselves (52.5\%)'' \citep{gkaintartzi_kiliari_tsokalidou2015,acquistapace2014}.\footnote{This attitude can be directly connected to an (often) subconscious process of linguicism ``used to construct some languages negatively as non-resources, even as handicaps, which may prevent children from enjoying the `full benefits' of majority languages'' \citep{nic_craith2006}. It may also be argued that teachers' formative process is itself embedded in monolingualism so that mainstream pedagogical approaches, adopted even by teachers who are themselves heritage language speakers, (un-)consciously favor a transparent language ideology. Indeed, majority language teachers have a ``certain perspective in regard to heritage language maintenance, and the importance of integrating into the majority culture and language'' (\citealt{weekly2020}, para. 25).}


Schools, therefore, often see the languages of new minorities as necessary but negative temporary tools while the child is learning the state's official language \citep{skutnabb-kangas2008}. As soon as he/she is deemed in some way competent in the language of instruction, the minority language can be left behind, and the child has no right to maintain this language and develop it further in the educational system, which is usually a one-way mechanism to assimilation.

The perceived burden of teachers to deal with bi- and multi-lingualism reflects an ongoing transformation into a multilingual society. And this process takes place in one of the functional parts of society, the educational system, which has to deal with contrary requests: a state-defined monolingualism and a multilingual reality \citep{decarlo2005}.

\textit{Policies} and \textit{practices} that countries use for language support can be divided into five different categories \citep{christensen_stanat2007}:

\begin{itemize}

\item Immersion: These programs provide no specific language support. students are ``immersed'' in the language of instruction within mainstream classrooms.

\item Immersion with systematic language support: Students are taught in the mainstream classroom, but they receive specified periods of instruction aimed at increasing proficiency in the language of instruction over a period of time.
\item Immersion with a preparatory phase: Students participate in a preparatory program before making the transition to mainstream classes.\footnote{Several countries have adopted programs to teach young children the language of instruction prior to attending compulsory education.  For instance, Germany has such programs for children who were born in the country or came to Germany at a very young age. The Flemish community in Belgium as well as Lithuania, Luxembourg, and Norway provide reception classes to children to equip them with language skills prior to attending compulsory education.  In addition, the Czech Republic, Finland, and some municipalities of Sweden offer pre-primary language instruction \citep{papademetriou2009}.}
\item Transitional bilingual schooling: Students initially learn in their minority language before teaching gradually shifts to the language of instruction.

\item Bilingual schooling \textit{strictu sensu}: Students receive significant amounts of instruction in their minority language; programs aim to develop proficiency both in the minority and the state official language.\footnote{Some Scandinavian countries, including Finland, Norway, and Sweden, as well as Cyprus, Estonia, and Latvia, offer bilingual instruction, where teachers teach in the new minorities' languages and the language of instruction \citep{papademetriou2009}. In the United  Kingdom, schools  have  always  been  able  to  offer  languages  spoken  by  their  pupils  within  the  modern  foreign languages curriculum, if they so wish. For instance, in the UK a number of schools are offering classes in Polish language and culture \citep{eurydice2009}.}

\end{itemize}


In Europe, most countries offer monolingual programs that provide additional support for second‐language learning. The most common approach in primary and secondary school is immersion with systematic language support \citep{christensen_stanat2007}. According to Christensen and Stanat, it is difficult to determine the extent to which the different language support programs contribute to the relative achievement levels of immigrant students \citep{christensen_stanat2007}. Studies conducted in this field indicate, however, that countries that tend to have long‐standing language support programs with clearly defined goals and standards are those with relatively small achievement gaps between new minorities and students from majority groups, or smaller gaps for second‐generation students compared to first‐generation students \citep{oecd2006}. In contrast, in countries where new minorities perform at significantly lower levels than their peers from the majority group, language support tends to be less systematic \citep{oecd2006}.

In terms of language support, EU legislation introduced in the 1970s following the Directive 1977/486/EEC requires MSs, as seen earlier, to provide supplementary language tuition for children of EU migrant workers, in both the language of the country of residence and in their minority language, with a view to facilitate their integration in the country of residence and in their country of origin should they subsequently return \citep{eu_coe_fra2015}. However, while this seems to offer quite generous and valuable supplementary support to children following their admission to a school in the settlement state, its implementation across different countries has been notoriously patchy and increasingly impractical given the range of different languages to accommodate \citep{eu_coe_fra2015}. Additionally, language education for new minorities frequently intersects with the teaching of foreign languages that are perceived to have high economic value, potentially posing challenges for minority students \citep{orourke_zhou2018}.

\subsection{Labor and Employment }\label{ch2.4.ss.4.2}

Also, in the field of labor and employment, language is considered to have an eminent role as an instrument to assure the successful integration of new minorities in society. Language is however also relevant to guarantee that the economy of a state functions properly. As for education, language has then a bidirectional dimension or value: a tool for new minorities to integrate and a tool for the economy to function properly.

Considerable research has examined the relationship between language proficiency in the state language and labor market integration. It is acknowledged that new minorities with greater language proficiency in the state language earn more and work in more skilled occupations than those with low proficiency, even after controlling for differences in education and skill associated with language abilities \citep{mchugh_challinor2011}. In particular, for the highly skilled, language fluency allows new minorities to practice the profession in which they are trained, rather than being downgraded to less-skilled jobs. Finally, for those with low- to mid-level skills, language proficiency is the ticket to better-paying jobs and upward mobility \citep{mchugh_challinor2011}.

Insufficient knowledge of the state language is in fact one of the main grounds for exclusion of new minorities in the labor market as this factor is so relevant for improving their status and to achieve social mobility. There is however evidence that language competence is often used as a pretext by many employers to exclude some candidates on cultural or religious grounds, in particular Muslims, even if they have been living in the country for 2 or 3 generations \citep{dechief_oreopoulos2012}.

Language-based discrimination in the labor market is a salient issue that links language policies and the equality principle. Language proficiency requirements and certification of language proficiency may indeed constitute a disproportionate obstacle for access to certain occupations for persons belonging to new minorities. In this regard, but only as EU nationals are concerned, the European Court of Justice acknowledged that it is legitimate to require a job applicant to have a certain level of linguistic knowledge. However, the Court stated: ``the right to require a certain level of knowledge of a language in view of the nature of the post must not encroach upon the free movement of workers. The requirements under measures intended to implement that right must not in any circumstances be disproportionate to the aim pursued and the manner in which they are applied must not bring about discrimination against nationals of other Member States''. \citep{ecj2015_belgium}.

Various European countries have rather strict state language laws and apply rigid requirements of state language proficiency to access positions in the public service and even jobs in the private sector. In practice, this may lead to direct or indirect discrimination against those who have insufficient command of the state language.

In parallel with this trend, the awareness of the importance of language skills, including languages of new minorities, is increasingly gaining ground in the current economic context. A multilingual and multicultural labor force is indeed considered crucial for economic growth and better jobs, enabling businesses to be more competitive \citep{huddleston_wolffhardt2016,npld2015,gazzola_wickstrom2016}. According to this view, languages, including languages of new minorities, are conceived as an opportunity for growth and prosperity.

\subsection{The Law Enforcement System}\label{ch2.4.ss.4.3}

The use of languages other than the state language in the law enforcement system has both an instrumental and an intrinsic value. Besides ensuring communication with the public authorities, using a language other than the state official in the public realm, shall guarantee, as emphasized earlier, the preservation and promotion of specific identities.

A number of safeguards in this field are set up by the European Convention on Human Rights (ECHR). Under Article 5 (2) ECHR, anyone arrested must be informed promptly and in a language that he/she understands of the reasons why he/she was arrested, so as to be able to seek judicial review of his/her arrest or detention.\footnote{In \textit{Zamir v. the United Kingdom}, the Commission was of the opinion that free legal aid should have been made available to an illegal immigrant detained pending deportation because of his poor command of English and the complexity of the case \citep{ECommHRa1983}.} Use of the terms `arrested `and `charge' in article 5(2) implies that the text belongs to a criminal law context \citep{schabas2015}. The information referred to in Article 5(2) need not be related in its entirety by the arresting officer at the very moment of the arrest. Besides, there is no requirement that the reasons be in writing or in any special form \citep{schabas2015}. This guarantee is applied in many cases of extradition and expulsion \citep{schabas2015}.

Human rights standards recognize also minimum guarantees in criminal prosecutions because it is clear that an accused person who is not familiar with the language used by the criminal tribunal is at a particular disadvantage. These guarantees include the right to receive information about the nature and cause of the charge in a language which he/she understands (art. 6 (3) (a) ECHR). This right is operative well before the trial and from the moment when the individual is `charged' \citep{schabas2015,ovey_white2006}.

The ECHR recognizes also the right of a person charged with a criminal offense to have free assistance by an interpreter if he/she cannot understand or speak the language used in court (art.  6(3)(e)). An evaluation of the linguistic abilities of the person charged with an offense is obviously vital to the application of this provision. Assessing the language skills of the accused person is primarily a matter for the domestic authorities, and, in principle, the ECHR will not intervene in such determinations. On occasion, however, the Court has made its own evaluation \citep{ecthr2010_katritsch}. The right to an interpreter applies not only to oral statements at the trial, but also to documentary material in order to have the benefit of a fair trial \citep{schabas2015,ovey_white2006}. The right to a free interpreter only extends to the language used in court: an accused who understands that language cannot insist upon the services of an interpreter to allow him/her to conduct his/her defence in another language, including a language of an ethnic minority of which he/she is a member \citep{ecommhr_zamir1983,ecommhr_bideault1986,ecthr2003_lagerblom}.

In addition to human rights norms --- the so-called `hard law' --- also `soft law' can be of assistance to new minorities as in the case of the recent Graz Recommendations on the Access to Justice for National Minorities \citep{osce_hcnm2017}. However, the decision on whether the scope of application of these recommendations should also be extended to new minorities is bound to be political and therefore it is largely left to the discretion of each country.

\section{Conclusions}\label{ch2.4.ss.5}

As international mobility flows continue to increase in Europe as elsewhere to an unprecedented high level, the question of integration through diversity governance reveals unequivocal urgency for most European countries that consider themselves to be rather culturally homogenous and cohesive. As a result, the process of integration of new minorities is seen as an important and urgent strategy to be adopted by most countries in order to retain an adequate level of social cohesion and prosperity \citep{collett_petrovic2014}.

In this chapter we have argued that in today's increasingly diverse societies it would be conceptually meaningful and beneficial in terms of diversity governance to consider widening the scope of minority rights traditionally conceived for old, historical minorities, such as those found in the CoE Framework Convention on National Minorities, to new minorities originating from migration. This would fill a still existing gap, especially in terms of rights related to identity and diversity, on which, we have seen, most international instruments on migrants' rights contain only weak and ambivalent references. But the protection of identity and language rights, including those of new minorities, is one of the bases of a veritable process of inclusion\footnote{In this regard, the OSCE HCNM called on States Parties to include minority cultures and tolerance in the general curriculum in order to increase inter-ethnic understanding and dialogue and to ``advance the important goal of sensitizing students to foreign cultures that do not qualify as national minorities, e.g., recent immigrants or refugees resident in the country'' \citep{osce_hcnm1997}. Likewise, article 6 FCMN calls on State Parties to encourage a ``spirit of tolerance and intercultural dialogue'' and promote ``mutual respect and understanding'', and the CSCE/OSCE Copenhagen Document invites State Parties to ``promote a climate of mutual respect, understanding, cooperation and solidarity among all persons living on its territory, without distinction as to ethnic or national origin or religion'' \citep{csce_osce1990}. See also \citet{eu2005_commonagenda,eu_ec_handbook_integration,eu_council_meeting2618_2004,osce_hcnm2006}. Among the vast literature on identity, see \citet{parekh2008,guild2004,nazroo_karlsen2003}.} through which minority groups can develop a genuine sense of loyalty and common belonging with the rest of the population without the threat of being forcibly assimilated in the mainstream society which, as a result, can engender resistance and alienation.\footnote{On the politics of belonging, see \citet{yuval-davis2006_art,geddes_favell1999}.}

Accommodating diversity, including linguistic diversity, is a powerful tool to reduce tensions and prevent conflicts; language is an important factor of diversity and identity building at individual and group level, thus some forms of accommodation should be found with all communities in society that speak languages different from the state language. These communities also include new minorities who are increasingly the most relevant diverse communities in contemporary societies.

In addition, cultural attachment to a language other than the official state language neither competes with nor replaces the requirement of skills and fluency in the official language of the country. A common public language is indeed necessary for the state to function and proficiency in the state language is an important tool by which integration is assured \citep{poggeschi2013}.

It is important to note, however, that in order to spark solidarity and social trust, and thus social cohesion, recognition of diversity is not sufficient: states should also tackle social structures and mechanisms that result in systematic deprivation and exclusion of new minorities from their equal public standing \citep{balint_latour2013}.

In conclusion, if a country is serious about the integration of new minorities, then it should not oppose the extension of the scope of application of minority provisions, including those pertaining to language rights to them. This would be an appropriate political gesture that underlines the importance of the country's integration policy and sends out a powerful message that populations of immigrants or asylum seekers are no longer seen as the ``legal Other'' \citep{acosta_arcarazo2011}\footnote{In a similar vein, Weiler notes: ``Concern for the cultural, linguistic and ethnic integration of aliens may send a signal which accentuates both the otherness of the alien and an intolerance of the dominant culture towards such otherness'' \citep{weiler1992}.} but as an integral, though distinct, part of the nation.

\section*{Abbreviations}
\begin{tabularx}{.49\textwidth}{@{}lQ@{}}
CoE & Council of Europe \\
EU & European Union \\
ACFC & Advisory Committee on the Framework Convention for the Protection of National Minorities \\
UN & United Nations \\
OSCE & Organizarion for Security and Cooperation in Europe \\
HCNM & High Commissioner on National Minorities \\
LTRD & EU Directive on the Status of Third-Country Nationals who are Long Term Residents \\
\end{tabularx}
\begin{tabularx}{.5\textwidth}{@{}lQ@{}}
MSs & Member States \\
UK & United Kingdom \\
PISA & Programme for International Student Assessment \\
OECD & Organization for Economic Co-operation and Development \\
NPLD & European Network to Promote Linguistic Diversity \\
ECHR & European Convention on Human Rights \\
ECtHR & European Court of Human Rights \\
\\
\\
\end{tabularx}

\section*{Acknowledgements}


This chapter draws on previous publication by \fullcite{marko_medda-windischer2018}.





\sloppy

\printbibliography[heading=subbibliography,notkeyword=this]

\end{document}

